\subsection{Derived Functors Preview}

Throughout this subsection, $R$ denotes a nonzero commutative ring with identity $1_R$. All modules are $R$-modules.

\begin{definition} \label{def:projective-injective-modules} \hlt{Projective and Injective Modules}\\
Let $P$ and $I$ be $R$-modules.
\begin{enumerate}[label=(\roman*)]
\setlength{\itemsep}{0pt}
\item The module $P$ is \hlt{projective} if each epimorphism $q:E\to N$ and homomorphism $f:P\to N$ admit $\widetilde f\in\operatorname{Hom}_R(P,E)$ with $q\widetilde f=f$.\\
The map $\widetilde f$ is a \hlt{lift} of $f$ through $q$.
\item The module $I$ is \hlt{injective} if each monomorphism $j:A\to B$ and homomorphism $f:A\to I$ admit $\widetilde f\in\operatorname{Hom}_R(B,I)$ with $\widetilde f j=f$.\\
The map $\widetilde f$ is an \hlt{extension} of $f$ along $j$.
\end{enumerate}
\end{definition}

\begin{theorem} \label{thm:basic-projective-constructions} \hlt{Basic Projective Constructions}\\
Assume the Choice Principle \ref{princ:axiom-of-choice-analysis}.
\begin{enumerate}[label=(\roman*)]
\setlength{\itemsep}{0pt}
\item Every free $R$-module is projective.
\item Every retract of a projective $R$-module is projective.
\item Every direct summand of a free $R$-module is projective.
\end{enumerate}
\begin{proof}
\begin{enumerate}[label=(\roman*)]
\setlength{\itemsep}{0pt}
\item Let $F$ be free with basis $\mathcal B$, let $q:E\to N$ be an epimorphism, and let $f:F\to N$ be a homomorphism.\\
Then $(\forall\,b\in\mathcal B)\,q^{-1}(\{f(b)\})\neq\emptyset$ by surjectivity of $q$. Choose $(e_b)_{b\in\mathcal B}\in\prod\limits_{b\in\mathcal B}q^{-1}(\{f(b)\})$ by the Choice Principle.\\
The assignment $b\mapsto e_b$ extends uniquely to $\widetilde f\in\operatorname{Hom}_R(F,E)$ by Theorem \ref{thm:free-module-map-extension}.\\
Maps $q\widetilde f$ and $f$ agree on $\mathcal B$, so $q\widetilde f=f$ by the uniqueness clause of that theorem; $F$ is projective by Definition \ref{def:projective-injective-modules}(i).
\item Let $Q$ be a retract of a projective module $P$, with $i:Q\to P$, $r:P\to Q$, and $ri=\operatorname{id}_Q$.\\
Let $q:E\to N$ be an epimorphism and let $f:Q\to N$ be a homomorphism.\\
Choose $h:P\to E$ with $qh=fr$ by projectivity of $P$; $q(hi)=f(ri)=f$, so $hi$ lifts $f$; $Q$ is projective by Definition \ref{def:projective-injective-modules}(i).
\item Let $Q$ be a direct summand of a free module $F$. Then $Q$ is a retract of $F$ by Corollary \ref{cor:retract-iff-direct-summand-module}.\\
The module $F$ is projective by (i), so $Q$ is projective by (ii).
\end{enumerate}
\hlt{Structural Commentary}: Basis lifts establish freeness, and retractions transfer lifting problems to a projective ambient module.\\
\hlt{Key Insight}: Projectivity is stable under retracts.
\end{proof}
\end{theorem}

\begin{proposition} \label{prop:injective-submodules-split} \hlt{Injective Submodules Split}\\
Let $I$ be injective and let $j:I\to E$ be a monomorphism.
\begin{enumerate}[label=(\roman*)]
\setlength{\itemsep}{0pt}
\item The map $j$ has a left inverse $r:E\to I$.
\item The submodule $\operatorname{im}j$ is a direct summand of $E$.
\end{enumerate}
\begin{proof}
\begin{enumerate}[label=(\roman*)]
\setlength{\itemsep}{0pt}
\item Choose $r:E\to I$ with $rj=\operatorname{id}_I$ by injectivity of $I$, applied to $\operatorname{id}_I:I\to I$.
\item We prove $E=\operatorname{im}j\oplus\ker r$. For $e\in E$, $e=j(r(e))+(e-j(r(e)))$ and $r(e-j(r(e)))=r(e)-(rj)(r(e))=0$.\\
Hence $E=\operatorname{im}j+\ker r$. If $x\in\operatorname{im}j\cap\ker r$, choose $a\in I$ with $x=j(a)$.\\
Then $a=(rj)(a)=r(x)=0$, so $x=0$. Thus the sum is direct.
\end{enumerate}
\hlt{Structural Commentary}: Extending the identity of $I$ produces a retraction, whose kernel is a complementary summand.\\
\hlt{Key Insight}: Injective modules split every monomorphism out of them.
\end{proof}
\end{proposition}

\begin{definition} \label{def:cochain-complex-cohomology} \hlt{Cochain Complexes and Cohomology}
\begin{enumerate}[label=(\roman*)]
\setlength{\itemsep}{0pt}
\item A \hlt{cochain complex} $K^\bullet=(K^n,\delta_K^n)_{n\in\Z}$ consists of homomorphisms $\delta_K^n:K^n\to K^{n+1}$ satisfying $(\forall\,n\in\Z)\,\delta_K^{n+1}\delta_K^n=0$.
\item The associated chain complex $K_\sharp$ is defined by $(K_\sharp)_m\defeq K^{-m}$ and $d_m^{K_\sharp}\defeq\delta_K^{-m}$.
\item Set $Z^n(K)\defeq\ker\delta_K^n$ and $B^n(K)\defeq\operatorname{im}\delta_K^{n-1}$. Then $B^n(K)\leq Z^n(K)$ by Proposition \ref{prop:chain-complex-image-kernel-containment} applied to $K_\sharp$.\\
The \hlt{$n$th cohomology module} is $H^n(K)\defeq Z^n(K)/B^n(K)=H_{-n}(K_\sharp)$.
\end{enumerate}
\end{definition}

\begin{definition} \label{def:projective-injective-resolutions} \hlt{Projective and Injective Resolutions}
\begin{enumerate}[label=(\roman*)]
\setlength{\itemsep}{0pt}
\item A \hlt{projective resolution} of $M$ consists of a chain complex $P_\bullet=(P_n,d_n)_{n\in\Z}$ and an epimorphism $\varepsilon:P_0\to M$.\\
It satisfies $(\forall\,n<0)\,P_n=0$, $\ker\varepsilon=\operatorname{im}d_1$, $(\forall\,n\in\N)\,\operatorname{im}d_{n+1}=\ker d_n$, and each $P_n$ for $n\in\N_0$ is projective.
\item An \hlt{injective resolution} of $N$ consists of a cochain complex $I^\bullet=(I^n,\delta_I^n)_{n\in\Z}$ and a monomorphism $\eta:N\to I^0$.\\
It satisfies $(\forall\,n<0)\,I^n=0$, $\operatorname{im}\eta=\ker\delta_I^0$, $(\forall\,n\in\N)\,\operatorname{im}\delta_I^{n-1}=\ker\delta_I^n$, and each $I^n$ for $n\in\N_0$ is injective.
\end{enumerate}
\end{definition}

\begin{proposition} \label{prop:hom-cochain-complex} \hlt{The Hom Cochain Complex}\\
Let $M$ be an $R$-module and let $I^\bullet=(I^n,\delta_I^n)_{n\in\Z}$ be a cochain complex. For $n\in\Z$, set $C^n\defeq\operatorname{Hom}_R(M,I^n)$ and $\Delta^n(\varphi)\defeq\delta_I^n\varphi$.
\begin{enumerate}[label=(\roman*)]
\setlength{\itemsep}{0pt}
\item Each $\Delta^n:C^n\to C^{n+1}$ is an $R$-module homomorphism.
\item The family $C^\bullet=(C^n,\Delta^n)_{n\in\Z}$ is a cochain complex.
\end{enumerate}
\begin{proof}
\begin{enumerate}[label=(\roman*)]
\setlength{\itemsep}{0pt}
\item Fix $(a,b)\in R^2$ and $(\varphi,\psi)\in(C^n)^2$. Then $\Delta^n(a\varphi+b\psi)=a\Delta^n(\varphi)+b\Delta^n(\psi)$ by $R$-linearity of $\delta_I^n$.
\item For $\varphi\in C^n$, $\Delta^{n+1}\Delta^n(\varphi)=\delta_I^{n+1}\delta_I^n\varphi=0$ by Definition \ref{def:cochain-complex-cohomology}(i). Hence $C^\bullet$ is a cochain complex by that definition.
\end{enumerate}
\hlt{Structural Commentary}: The covariant variable of $\operatorname{Hom}_R(M,-)$ preserves the direction of the coboundary maps.\\
\hlt{Key Insight}: Applying $\operatorname{Hom}_R(M,-)$ to an injective resolution produces a cochain complex.
\end{proof}
\end{proposition}

\begin{definition} \label{def:chosen-resolution-ext} \hlt{Chosen-Resolution Model for Ext}\\
Let $M,N$ be $R$-modules, and let $(I^\bullet,\eta)$ be an injective resolution of $N$. For $n\in\N_0$, set $\operatorname{Ext}_R^n(M,N;I^\bullet)\defeq H^n\bigl(\operatorname{Hom}_R(M,I^\bullet)\bigr)$.\\
The resolution remains in the notation because resolution independence is not part of this definition.
\end{definition}

\begin{perspective} \label{persp:tor-ext-resolution-models} \hlt{Resolution Models for Tor and Ext}\\
After tensor products of arbitrary $R$-modules and their induced maps are developed, a projective resolution $P_\bullet\to M$ produces a tensor complex with terms $P_n\otimes_RN$. Its homology is the projective-resolution model for $\operatorname{Tor}_n^R(M,N)$, parallel to Definition \ref{def:chosen-resolution-ext}.\\
A comparison theorem is still required to identify the models from different resolutions by canonical isomorphisms.\\
Until that theorem and the module-level tensor construction are available, the abbreviations $\operatorname{Tor}_n^R(M,N)$ and $\operatorname{Ext}_R^n(M,N)$ are forward notation rather than proof authority.
\end{perspective}
